\chapter{Computation, Representation, and the Chaitin Barrier}
\label{chap:computation-chaitin}

The construction can now count, order, approximate the continuum, and keep time. The natural next
question is what it can \emph{compute} --- and the honest answer comes with a hard edge. There is
exactly one barrier to computation known to be physical rather than merely practical, and the
construction walks straight up to it and calibrates against it. This chapter is about that edge: what
computation is in a finite world, where the symbols it manipulates actually come from, why noise and
signal are produced by the very same process, and how, at the bottom, the resolution of any
instrument runs out. Along the way the construction takes its first real step toward calculus, by
noticing that an observation is a velocity and the area under it is an integral.

\section{Computation and representation}
\label{se:computation-and-representation}

A computation, in the construction, is not a magic box. It is a sequential refinement of a single
record: a process that reads the record, appends one more determined step, and repeats, traversing its
own internal structure in a controlled order. This is exactly the content of the experiment the
construction files under Turing's name --- that any finite-dimensional process can be represented as
the orderly refinement of one record, provided the instrument can decompose its structure and walk
through it in a coordinated way. There is nothing here but the disciplined extension of a ledger, one
step at a time, which is all a Turing machine ever was.

What the construction is scrupulously honest about is that it gets no guardrails. To run a computation
it must already know that the computation halts; nothing inside the process warns it if the process
will run forever. The halting problem is not an external theorem imported for color --- it is the
ambient condition under which every computation here operates, the standing reminder that a finite
device cannot be promised in advance that its own procedure will terminate. The construction proceeds
anyway, on processes it has reason to believe halt, and is candid that the belief is a precondition,
not a guarantee.

There is a deeper honesty in this section, and it is the one most worth carrying forward. Where do the
symbols a computation manipulates come from? Not from mathematics. The construction states this
flatly: a representation is not given by the mathematics; the symbols come from elsewhere. A choice of
how to write a quantity --- which marks stand for which distinctions --- is imposed from outside the
formal system, by a community that agrees a pattern in one head is congruent to a pattern in another
and in the devices that carry it between minds. Mathematics can manipulate a representation once it has one, but it
cannot hand itself the representation to begin with; that act is prior to the formal system and
external to it. This is not a defect. It is the seam where a formal construction meets the world that
uses it, and naming the seam honestly is what lets the later chapters attach physical words to formal
objects without pretending the words were derived.

A small illustration fixes the point. The same quantity --- three --- can be written as a row of three
tally marks, as a digit string in base ten, as a different digit string in base two, as a position on
a dial, as a length of rod, or as a pattern of voltages on a wire. Mathematics has no opinion about
which of these is the number; they are all equally faithful, and the relations among numbers hold
identically whichever is chosen. The choice between them is made for reasons that live entirely outside
the mathematics --- legibility to a human eye, cheapness to a circuit, robustness against noise on a
line --- by whoever has to read the result. This is what it means to say representation is not given:
the formal system supplies the structure, the world supplies the symbols, and the two meet only by an
agreement the formal system cannot itself certify. Every later chapter that pins a physical name to a
constructed object spends exactly this kind of agreement, and it is cleaner to see the currency now.

\section{Calibrating against Chaitin}
\label{se:calibrating-against-chaitin}

To find the edge of computation you need something to push against, and the construction chooses the
hardest object available: Chaitin's constant, the probability that a random program halts. It is a
genuinely strange number --- in one sense completely determined, a single real number fixed by the
machinery, and in another sense completely unpredictable, its digits encoding the answer to the
halting problem and so computable by no finite procedure. The construction does not try to compute it,
which would be hopeless. It uses it as a calibration target: the one universally definable constant
whose behavior is perfectly lawful and perfectly unforeseeable at once, the standing physical barrier
that any honest account of computation must run into rather than around.

The reason this matters for measurement is the chapter's sharpest point, and it is the lesson of the
experiment under Chaitin's name. A record may consist entirely of finite, distinguishable, perfectly
admissible events and still admit no extractable law whatsoever. Piling up facts does not guarantee
that a truth will emerge from them; some finite records are simply lawless, noise all the way down,
and no amount of further refinement will distill a rule from them because there is no rule there to
distill. And here is the unsettling part: noise of this kind is produced by the \emph{same} process
that produces signal. The construction generates a candidate sequence the only way it can, by
refinement, and whether the result is meaningful signal or pure noise is not something the generating
process can tell you. Signal and noise are the same act with different luck --- the construction \emph{funges} them into one
process, the two set apart only by which way the luck fell.

This puts the construction in the position of a gambler, and it is candid about the odds. Having
generated a sequence by refinement, all it can do is bet that the sequence is lawful and watch the bet
ride. Each further bit that comes out as predicted is a small payout, a little more evidence that a
rule is present; each bit is equally a fresh chance to be wrong. There is no point at which the
winnings convert into a certificate, because the very object it calibrates against --- Chaitin's
constant --- is the standing proof that a perfectly definite sequence can be perfectly unforecastable.
The construction does not find this discouraging; it finds it clarifying. Measurement has always been
the accumulation of confirmations against the permanent possibility of refutation, and computation,
pushed to its physical edge, simply makes the asymmetry impossible to look away from.

\begin{quote}
\emph{A worked check.} Suppose you and the construction each generate a stream of bits meant to be
the same constant, computed independently. Compare them in order. As long as they agree, bit after
bit, your confidence that you are both computing the same lawful object climbs --- this is the study
of the previous chapter, funging the agreeing trials and rounding toward ``true''. But the agreement is
never a proof. A single disagreement at the thousandth bit, or a permutation that reveals you were
computing the same bits in a different order, collapses the verdict, and nothing in the stream warned
you it was coming. You cannot certify in advance that the stream is signal; you can only watch it fail
to be noise, for now. That asymmetry --- endless tentative confirmation, instant refutation --- is
Chaitin's barrier made into a daily experience, and it is exactly Hume's induction from the last
chapter wearing a computer's clothes.
\end{quote}

\section{The machine epsilon}
\label{se:the-machine-epsilon}

Every finite instrument has a smallest step it can resolve, a floor below which two readings are the
same reading. The construction makes this floor explicit and calls it, in effect, its machine epsilon:
the least distinguishable increment, the size of the smallest move the apparatus can certify as a
move at all. Below the machine epsilon there is no signal and no noise, only sameness; it is the scale
at which the Berkeley--Galileo lesson of the first chapter bites, where two descriptions that differ
by less than epsilon are, for this instrument, one description.

With a resolution floor in hand, the construction can compare two quantities the way a finite device
actually does: by bisection. To decide which of two numbers is larger, or to locate a value, it
repeatedly tanges the half that must contain the answer --- selecting, by a single decidable question,
which side of a split the target lies on --- and discards the other half, halving the remaining
uncertainty at each step. The procedure does not run forever; it runs until the surviving interval is
narrower than the machine epsilon, at which point the remaining ambiguity is below resolution and the
comparison simply stops with a verdict. Comparison, in a finite world, is not an act of seeing which
number is bigger; it is a finite sequence of decidable cuts that terminates when further cutting would
fall beneath the instrument's own floor.

Made concrete, the procedure is the one every instrument with a dial secretly runs. To pin a value
known to lie between zero and one, ask first whether it exceeds one half; the answer, a single
decidable bit, tanges the surviving half. Ask next whether it exceeds the midpoint of that half, and
again, and again, each question chopping the remaining interval in two and each answer a clean yes or
no. After ten such questions the interval is a thousandth of its original width; after twenty, a
millionth. The construction does not press on forever: it stops the first time the surviving interval
is narrower than its machine epsilon, because at that point the two ends of the interval are the same
reading and no further question could separate them. The value has not been found exactly --- exactly
is not available to a finite device --- but it has been cornered to within the instrument's resolution,
which is all a measurement ever claims.

\section{From observation to area}
\label{se:from-observation-to-area}

The chapter's last move is the one that opens the door to everything calculus-shaped that follows, and
it begins by noticing what an observation actually is. The construction models the simplest possible
physical event --- the moment a thing at rest begins to move --- and reads the observation off as a
\emph{velocity}: zero while the object is held, one once it slips. An observation is not a static
quantity but a rate, the first derivative of position dressed as a measurement. The event it watches
is the onset of motion under static friction, and the construction chooses that example for a precise
reason: it is the cleanest case of something a finite ledger cannot predict.

This is the experiment under the names of Da Vinci and Coulomb. The force needed to hold a thing
still rises with the load by Coulomb's simple law, but the exact instant at which static friction
gives way and the object slips is a finite, distinguishable event that cannot be read off from any
amount of prior refinement of the ledger. You can press harder and harder, refining the record
indefinitely, and the record will not tell you when the slip will come; only the slip itself, when it
happens, settles it. This unpredictable onset is what the construction means by a curvature term: the
genuinely nonlinear part of a process, the part that cannot be reached by extrapolating the smooth
trend, the place where the second-order structure asserts itself. Friction is the construction's first
curvature, and it earns the name by being exactly the thing linear prediction misses.

Once an observation is a velocity, the quantity that accumulates it is an area, and the construction
builds that next: the area under a sequence of observations, assembled the same way numbers were, one
strip at a time. That area is a first integral --- velocity accumulated into displacement, the
integral that the calculus of the coming chapters will build on, and the place the construction first
admits that it is doing integration rather than mere counting. It is worth noting that area arrives
with the opposite orientation from the numbers of the first chapter: where those accumulated most
significant part first, the integral accumulates from the other end, the contravariant counterpart to
the covariant tower --- the two ways of building meeting again, now as derivative and integral.

It is worth saying why the construction reaches derivative and integral here, two chapters before it
strictly needs them. A velocity is what a measurement reports when it watches a quantity change; an
area is what accumulates when those reports are summed back up. The two are inverse operations, and
between them they let the construction ask the question the whole middle of the book turns on: not what
a quantity \emph{is} at an instant, but how it \emph{responds} when its surroundings are varied, and
what survives unchanged when they are. Reading observation as a rate and area as its integral is the
minimal equipment for that question --- a way to register change and a way to total it --- and having
built both from finite strips, the construction can treat a measured history as something with a slope
and an area rather than a bare list of numbers. That is the threshold of the calculus of variations,
and the friction example was the construction stepping up to it: a curvature term is precisely what a
theory of how things vary must be able to see.

A coda to the friction example guards against the wrong conclusion. The experiment under Laplace's
name observes that, relative to a single committed record, each next refinement looks uniquely
determined by its predecessors --- the record appears to run like Laplace's demon, the future fixed by
the past. The construction keeps the appearance and demotes the demon. Determinism here is a property
of a finite record's internal consistency, true relative to the anchors already committed, not a
cosmic guarantee that the next observation cannot surprise. The demon is real but local, and the
Chaitin barrier is the reminder that its writ does not extend past the edge of what has actually been
recorded.

\section{What is demonstrated, and what is assumed}
\label{se:demonstrated-assumed-ch4}

Keeping to the experimental library's discipline, we separate the built from the named. The chapter
reaches toward large ideas --- computation, determinism, the calculus --- and the modest truth
underneath them is worth isolating.

What is demonstrated is a finite theory of computation and its first calculus. There is a notion of
computation as the sequential refinement of one record, with the halting problem honestly marked as
the condition it runs under. There is a calibration object, Chaitin's constant, used not by being
computed but by being pushed against. There is a machine epsilon, a least resolvable step, and a
bisection comparison that tanges its way to a verdict and halts at that floor. And there is an
observation read as a velocity and an area built as its first integral, both finite, both computable,
both checkable. None of this assumes infinite precision, an oracle for halting, or a completed
continuum.

\begin{quote}
\emph{A note on the labels.} Several names in this chapter are bridges, not theorems. Calling the
unpredictable onset of slip a \emph{curvature term}, an observation a \emph{velocity}, and an
accumulated area a \emph{first integral} are modeling choices that align the finite constructions with
the vocabulary of physics and calculus; the constructions are earned, the physical readings are laid
over them. The deepest assumption in the chapter is the one the construction states outright: that a
\emph{representation} can be had at all. The symbols come from outside the mathematics, and the whole
edifice runs on the prior, unproven agreement that a pattern here stands for a pattern there. That
agreement is the price of admission, paid before the first line of construction, and honesty requires
naming it as an assumption rather than a result.
\end{quote}

What is assumed, then, is a representation to compute with and a belief that the procedures run will
halt --- the two things a finite computer cannot give itself and must be handed. With those granted,
the construction has computation, a resolution floor, and the first integral of a velocity, which is
precisely the equipment needed to take the next step: to treat a measurement not as a single reading
but as something that varies, and to ask how it responds when its surroundings are changed. That is
the move into gauge, superposition, and the local present, and it is where the construction stops
describing arithmetic and starts describing fields.
